\title{Chain-of-Thought Prompting Elicits Reasoning in Large Language Models}

\begin{abstract}
We investigate how producing a \textit{chain of thought}—a sequence of intermediate reasoning steps—substantially enhances the capability of large language models to tackle complex reasoning problems. Specifically, we demonstrate that such reasoning skills naturally arise in sufficiently large language models through a straightforward approach called \textit{chain-of-thought prompting}, which involves presenting a few chain-of-thought examples as prompts.

Tests conducted on three large language models reveal that chain-of-thought prompting boosts performance across various arithmetic, commonsense, and symbolic reasoning tasks. The improvements can be remarkable; for example, using just eight chain-of-thought exemplars to prompt a PaLM 540B model achieves state-of-the-art results on the GSM8K dataset of math word problems, even outperforming a finetuned GPT-3 equipped with a verifier.
\end{abstract}

\section{Introduction}
The field of NLP has recently undergone a transformation driven by language models \citep[][\textit{inter alia}]{peters-etal-2018-deep,devlin-etal-2019-bert,brown2020language}. Increasing the size of language models has been shown to bring various advantages, including better performance and enhanced sample efficiency \citep[\textit{inter alia}]{kaplan2020scaling,brown2020language}. Nonetheless, merely enlarging model size has not sufficed to achieve strong performance on difficult tasks such as arithmetic, commonsense, and symbolic reasoning \citep{rae2021scaling}.

In this work, we examine how a simple technique can unlock the reasoning potential of large language models, inspired by two main insights. First, methods for arithmetic reasoning benefit from generating natural language explanations that guide the derivation of the final answer. Earlier research has enabled models to produce natural language intermediate steps either by training from scratch \citep{ling-etal-2017-program} or by finetuning pretrained models \citep{cobbe2021training}, alongside neuro-symbolic approaches that utilize formal languages instead of natural language \citep{roy-roth-2015-solving, chiang-chen-2019-semantically,amini-etal-2019-mathqa, chen2019neural}. Second, large language models present an exciting opportunity for in-context few-shot learning through \emph{prompting}. Rather than finetuning a separate model for each new task, one can simply ``prompt'' the model with a handful of input-output examples illustrating the task. Remarkably, this method has been effective for a variety of straightforward question-answering problems \citep{brown2020language}.

Both of the aforementioned approaches, however, have significant drawbacks. For rationale-augmented training and fine-tuning techniques, generating a substantial collection of high-quality rationales is expensive and considerably more complex than using straightforward input--output pairs typical in standard machine learning. In contrast, the conventional few-shot prompting method employed by \citet{brown2020language} tends to perform inadequately on tasks that demand reasoning skills, and often does not show significant improvement with increasing language model size \citep{rae2021scaling}. In this study, we integrate the advantages of these two strategies in a manner that circumvents their respective limitations. Specifically, we investigate the capability of language models to execute few-shot prompting for reasoning tasks using prompts composed of triples: $\langle$input, \emph{chain of thought}, output$\rangle$.  
A \emph{chain of thought} denotes a sequence of intermediate natural language reasoning steps that culminate in the final output; we label this method as \textit{chain-of-thought prompting}. An example of such a prompt is provided in \cref{fig:pull-figure}.

We conduct empirical analyses across benchmarks involving arithmetic, commonsense, and symbolic reasoning, demonstrating that chain-of-thought prompting surpasses standard prompting, occasionally by a substantial margin. \cref{fig:pull-bar-chart} presents one such finding—on the GSM8K dataset of math word problems \citep{cobbe2021training}, chain-of-thought prompting with \palm{} 540B significantly outperforms conventional prompting and establishes a new state-of-the-art. A prompting-only approach is valuable because it does not necessitate a large training dataset and allows a single model checkpoint to address multiple tasks without loss of generality. This research highlights how large language models can learn effectively from a few examples with natural language task descriptions (as opposed to learning patterns solely from extensive training datasets).

\section{Chain-of-Thought Prompting} When tackling a complex reasoning problem, such as a multi-step math word problem, it is common to mentally break down the problem into intermediate steps and solve each one before arriving at the final answer: \textit{``After Jane gives 2 flowers to her mom she has 10 $\ldots$ then after she gives 3 to her dad she will have 7 $\ldots$ so the answer is 7.''} This paper aims to equip language models with the capability to produce a comparable \textit{chain of thought}—a logically connected sequence of intermediate reasoning steps that culminate in the final solution to a problem. We demonstrate that language models of sufficient size can generate chains of thought if few-shot prompting includes exemplars that showcase chain-of-thought reasoning.

\cref{fig:pull-figure} illustrates an example where a model generates a chain of thought to solve a math word problem it would otherwise answer incorrectly. In this instance, the chain of thought functions similarly to a solution and can be interpreted as one, though we prefer to term it a chain of thought to emphasize that it replicates a stepwise cognitive process toward the answer (also, solutions or explanations usually follow \textit{after} the final answer \citep[][\textit{inter alia}]{narang2020wt5,wiegreffe2021reframing,lampinen2022can}).

Chain-of-thought prompting offers several appealing features as a method for enhancing reasoning capabilities in language models.

\begin{enumerate}[topsep=1pt,itemsep=0ex]
    \item Firstly, chain of thought enables models to break down multi-step problems into intermediate components, allowing extra computational resources to be dedicated to problems that involve more complex reasoning steps.
    \item Secondly, chain of thought offers an interpretable insight into the model’s behavior, revealing the potential rationale behind a specific answer and facilitating the identification of errors in the reasoning process (although completely understanding the computations supporting an answer remains an unresolved challenge).
    \item Thirdly, chain-of-thought reasoning is applicable to tasks such as math word problems, commonsense reasoning, and symbolic manipulation, and is potentially extendable (at least theoretically) to any problem that humans solve using language.
    \item Lastly, chain-of-thought reasoning can be easily elicited from sufficiently large pre-trained language models by incorporating chain-of-thought examples into the few-shot prompt exemplars.
\end{enumerate}

In empirical studies, we will demonstrate the effectiveness of chain-of-thought prompting in the domains of arithmetic reasoning (\cref{sec:arithmetic-reasoning}), commonsense reasoning (\cref{sec:commonsense-reasoning}), and symbolic reasoning (\cref{sec:symbolic-reasoning}).

\section{Arithmetic Reasoning}\label{sec:arithmetic-reasoning}
We start by examining math word problems as illustrated in \cref{fig:pull-figure}, which test the arithmetic reasoning capabilities of language models. While these problems are straightforward for humans, language models frequently find arithmetic reasoning challenging \citep[][\textit{inter alia}]{hendrycks2021measuring,patel-etal-2021-nlp}. Remarkably, when combined with chain-of-thought prompting, the 540B parameter language model performs on par with specialized finetuned models across several benchmarks, even setting a new state-of-the-art on the demanding GSM8K dataset \citep{cobbe2021training}.

\subsection{Experimental Setup}

We investigate the application of chain-of-thought prompting on different language models across multiple benchmarks.

\textbf{Benchmarks.}
We evaluate five math word problem benchmarks:
\textbf{(1)} the \textbf{GSM8K} math word problems benchmark \citep{cobbe2021training},
\textbf{(2)} the \textbf{SVAMP} dataset containing math problems with varied structures \citep{patel-etal-2021-nlp},
\textbf{(3)} the \textbf{ASDiv} dataset covering diverse math word problems \citep{miao-etal-2020-diverse},
\textbf{(4)} the \textbf{AQuA} dataset focused on algebraic word problems, and
\textbf{(5)} the \textbf{MAWPS} benchmark \citep{koncel-kedziorski-etal-2016-mawps}. Example problems are provided in Appendix \cref{tab:math-datasets}.

\textbf{Standard prompting.} As a baseline, we utilize the standard few-shot prompting approach popularized by \citet{brown2020language}, where the language model is presented with a few input-output examples in context before generating an answer for a test instance. These exemplars are formatted as question-and-answer pairs, and the model directly outputs the final answer, as depicted in \cref{fig:pull-figure} (left).

\textbf{Chain-of-thought prompting.} Our method involves enhancing each exemplar in few-shot prompting with a chain of thought that corresponds to the given answer, as depicted in \cref{fig:pull-figure} (right). Since most datasets only provide an evaluation split, we manually created a set of eight few-shot exemplars that include chains of thought for use in prompting—\cref{fig:pull-figure} (right) illustrates one such exemplar, and the complete collection is presented in Appendix \cref{tab:appendix-mwp-prompts}. (These specific exemplars were not subject to prompt engineering; the robustness of the approach is analyzed in \cref{subsec:robustness} and \cref{subsec:faq-prompt-engineering}.) To evaluate whether this style of chain-of-thought prompting can effectively induce accurate reasoning across various math word problems, we applied this single set of eight chain-of-thought exemplars to all benchmarks except AQuA, which is formatted as multiple choice rather than free response. For AQuA, we selected four exemplars and their solutions from the training set, as detailed in Appendix \cref{tab:appendix-aqua-prompt}.

\textbf{Language models.} We assess five large language models. The first is \textbf{GPT-3} \citep{brown2020language}, employing text-ada-001, text-babbage-001, text-curie-001, and text-davinci-002 variants, which are likely to correspond to InstructGPT models with 350M, 1.3B, 6.7B, and 175B parameters respectively \citep{ouyang2022training}. The second model is \textbf{\lamda{}} \citep{thoppilan2022lamda}, available in sizes of 422M, 2B, 8B, 68B, and 137B parameters. The third is \textbf{\palm{}}, with 8B, 62B, and 540B parameter versions. The fourth is \textbf{UL2 20B} \citep{tay2022unifying}, and the fifth is \textbf{Codex} \citep[][code-davinci-002 in the OpenAI API]{chen2021evaluating}. Model outputs were generated using greedy decoding (although subsequent research demonstrates that chain-of-thought prompting performance can be enhanced by taking the majority vote of final answers across multiple sampled generations \citep{wang2022self}). For \lamda{}, we report the average performance over five random seeds, each corresponding to a different random permutation of the exemplars. Since \lamda{} experiments exhibited limited variance across seeds, to conserve computational resources we present results using a single exemplar ordering for all other models.

\subsection{Results}  
The most significant findings regarding chain-of-thought prompting are summarized in \cref{fig:main-math}, with comprehensive experimental results for each model group, size, and benchmark detailed in \cref{tab:all-lm-math} in the Appendix. There are three primary insights. First, \cref{fig:main-math} demonstrates that chain-of-thought prompting emerges as a capability with increasing model scale \citep{wei2022emergent}. Specifically, chain-of-thought prompting does not improve performance in smaller models and only produces performance enhancements when applied to models approximately 100B parameters or larger. Qualitative analysis showed that smaller models generated fluent but logically flawed chains of thought, resulting in worse performance compared to standard prompting.

\input{main_fables/main-math}  
Second, chain-of-thought prompting yields greater performance improvements on more complex problems. For example, on GSM8K (the dataset with the lowest baseline accuracy), performance more than doubled for the largest GPT and \palm{} models. Conversely, on SingleOp—the simplest subset of MAWPS requiring only one step for a solution—the performance gains were either negligible or negative (refer to Appendix \cref{tab:mawps-subsets}).

Third, chain-of-thought prompting applied to GPT-3 175B and \palm{} 540B models shows competitive results compared to previous state-of-the-art approaches, which typically involve fine-tuning a task-specific model on a labeled training set.  
\cref{fig:main-math} illustrates how \palm{} 540B leverages chain-of-thought prompting to establish new state-of-the-art performance on GSM8K, SVAMP, and MAWPS (noting that standard prompting already surpassed the prior best on SVAMP). For the other datasets, AQuA and ASDiv, \palm{} with chain-of-thought prompting achieves accuracy within 2\% of the state of the art (see Appendix \cref{tab:all-lm-math}).

To gain a deeper understanding of why chain-of-thought prompting is effective, we conducted a manual inspection of model-generated chains of thought produced by \lamda{} 137B on the GSM8K dataset. Among 50 randomly selected examples where the model provided the correct final answer, all but two of the generated chains of thought were both logically and mathematically accurate; the two exceptions happened to yield the correct answer by chance (refer to \cref{subsec:correct-chain-of-thought} and \cref{tab:appendix-gsm8k-correct-examples} for examples of correctly generated chains of thought). Additionally, we randomly reviewed 50 cases where the model's final answer was incorrect. Our analysis revealed that 46\% of these chains of thought were nearly correct but contained minor errors (such as calculation mistakes, symbol mapping errors, or a missing reasoning step), while the remaining 54\% exhibited significant flaws in semantic comprehension or coherence (see \cref{subsec:incorrect-chain-of-thought-analysis}).

To offer some insight into why scaling enhances chain-of-thought reasoning capabilities, we performed a similar error analysis comparing \palm{} 62B and whether those errors were resolved when scaling up to \palm{} 540B. The findings indicate that increasing the model size to 540B corrects a substantial fraction of the one-step omission and semantic understanding errors present in the 62B model (see \cref{subsec:why-scale-helps}).

\subsection{Ablation Study}\label{subsec:math-ablation}

The demonstrated advantages of chain-of-thought prompting naturally prompt the question of whether other forms of prompting can yield similar performance gains. \cref{fig:bar_ablation} presents an ablation study exploring three variants of chain-of-thought prompting, described below.

\textbf{Equation only.} One hypothesis for the effectiveness of chain-of-thought prompting is that it helps by generating the mathematical equation to solve the problem. To test this, we examine a variant where the model is prompted to produce only the mathematical equation before providing the final answer. As shown in \cref{fig:bar_ablation}, prompting for the equation alone does not significantly improve performance on GSM8K. This suggests that the semantic complexity of GSM8K questions is too high to be directly converted into equations without the intermediate natural language reasoning steps present in chain-of-thought prompting. However, for datasets involving one- or two-step problems, we observe that equation-only prompting does enhance performance, since the equations can be straightforwardly derived from the questions (see Appendix \cref{tab:ablations-arithmetic}).

\input{main_fables/ablation-bar}
\textbf{Variable compute only.} One hypothesis is that the chain of thought enables the model to allocate more computation (i.e., intermediate tokens) for more difficult problems. To separate the effect of variable computation from chain-of-thought reasoning, we experiment with a setup where the model is instructed to output only a sequence of dots ($\ldots$), with the number of dots matching the number of characters in the equation necessary to solve the problem. This variant achieves performance roughly equivalent to the baseline, indicating that variable computation alone does not explain the effectiveness of chain-of-thought prompting, and that there seems to be added value in articulating intermediate steps in natural language.

\textbf{Chain of thought after answer.} Another possible advantage of chain-of-thought prompting could be that such prompts help the model better retrieve relevant knowledge gained during pretraining. To examine this, we test an alternative configuration where the chain of thought prompt is provided only after the answer, to determine whether the model relies on the generated chain of thought when producing the final answer. This variant performs similarly to the baseline, suggesting that the step-by-step reasoning encapsulated in the chain of thought offers benefits beyond simply activating stored knowledge.

\input{main_fables/main-robustness}
\subsection{Robustness of Chain of Thought}\label{subsec:robustness}

The sensitivity to exemplar choice is an important aspect of prompting methods—for example, changing the order of few-shot exemplars can cause GPT-3’s accuracy on SST-2 to vary dramatically from near random chance (54.3\%) to near state-of-the-art performance (93.4\%) \citep{zhao2021calibrate}. In this final subsection, we assess the robustness of chain-of-thought prompting with chains authored by different annotators. Beyond the previous results that used chains written by Annotator A, two other co-authors (Annotators B and C) independently composed chains of thought for the same few-shot exemplars (see \cref{sec:appendix-alternate-annotators}).
Annotator A also produced another, more concise chain of thought compared to the original, modeled after the solution style in \citet{cobbe2021training}.\footnote{For example, while the original chain of thought consists of multiple short sentences (e.g., \textit{``There were originally 9 computers. For each of 4 days, 5 more computers were added. So 5 * 4 = 20 computers were added. 9 + 20 is 29.''}), the concise chain would read \textit{``5 * 4 = 20 new computers were added. So there are 9 + 20 = 29 new computers in the server room now.''}}

\cref{fig:bar-robustness-analysis} presents these findings for \lamda{} 137B on GSM8K and MAWPS (additional ablation results for other datasets can be found in Appendix \cref{tab:ablations-arithmetic} / \cref{tab:ablations-commonsense-symbolic}). While variability exists across different chain of thought annotations, as anticipated when employing exemplar-based prompting \citep{le-scao-rush-2021-many,reynolds2021prompt,zhao2021calibrate}, every set of chain of thought prompts significantly surpasses the standard baseline. This outcome suggests that effective chain of thought usage does not rely on any specific linguistic style.

To verify that chain-of-thought prompting is effective with other exemplar sets, we conducted experiments using three sets of eight exemplars randomly drawn from the GSM8K training data, an independent source (the examples in this dataset already contain reasoning steps resembling a chain of thought).\footnote{We selected examples with $\leq 60$ tokens to fit our input context window, and restricted them to $\leq 2$ steps to ensure a fair comparison with the eight exemplars we composed.} \cref{fig:bar-robustness-analysis} demonstrates that these prompts performed similarly to our manually crafted exemplars, also significantly outperforming standard prompting.

Beyond robustness to annotators, independently constructed chains of thought, different exemplars, and various language models, we also observe that chain-of-thought prompting for arithmetic reasoning remains stable across different orders of exemplars and varying numbers of exemplars (refer to \cref{subsec:faq-prompt-engineering}).

\section{Commonsense Reasoning}\label{sec:commonsense-reasoning}

Although chain of thought is especially well-suited for math word problems, its language-based nature renders it applicable to a wide range of commonsense reasoning tasks, which require reasoning about physical and social interactions assuming general background knowledge. Commonsense reasoning is essential for engaging with the world and remains beyond the capabilities of current natural language understanding systems \citep{talmor2022commonsenseqa}.

\textbf{Benchmarks.}  
We analyze five datasets that encompass a variety of commonsense reasoning tasks. The widely-used \textbf{CSQA} \citep{talmor-etal-2019-commonsenseqa} presents commonsense questions about the world, involving intricate semantics that often depend on background knowledge.  
\textbf{StrategyQA} \citep{geva-etal-2021-aristotle} challenges models to deduce multi-hop strategies for answering questions. We select two specialized evaluation sets from the BIG-bench initiative \citep{bigbench}: \textbf{Date} Understanding, which requires inferring a date based on provided context, and \textbf{Sports} Understanding, which involves assessing whether a sentence related to sports is believable or not. Lastly, the \textbf{SayCan} dataset \citep{ahn2022can} focuses on translating natural language instructions into sequences of robot actions from a predefined discrete set.  
\cref{fig:dataset-examples} presents examples with chain of thought annotations across all datasets.

\textbf{Prompts.}  
We adopt the same experimental protocol as described in the previous section. For CSQA and StrategyQA, we randomly picked examples from the training sets and manually constructed chains of thought to serve as few-shot exemplars. Since the two BIG-bench tasks lack training sets, we used the first ten examples from the evaluation set as few-shot exemplars and report results on the remaining examples. For SayCan, we utilized six examples from the training set used by \citet{ahn2022can}, manually crafting chains of thought for them as well.

\textbf{Results.}
\cref{fig:commonsense-results} presents these findings for \palm{} (comprehensive results for \lamda{}, GPT-3, and various model sizes are provided in \cref{tab:all-lm-commonsense}).
For every task, increasing the model size enhanced the performance under standard prompting; chain-of-thought prompting brought about additional improvements, with the most notable gains observed for \palm{} 540B. When using chain-of-thought prompting, \palm{} 540B demonstrated strong performance compared to baselines, surpassing the previous state of the art on StrategyQA (75.6\% vs 69.4\%) and outperforming an unaided sports fan in sports understanding (95.4\% vs 84\%).
These outcomes illustrate that chain-of-thought prompting can also boost performance on tasks that require diverse commonsense reasoning skills (although the improvement was minimal for CSQA).

\input{main_fables/main-commonsense}

\input{main_fables/main-symbolic}
\section{Symbolic Reasoning}\label{sec:symbolic-reasoning}
Our last experimental analysis focuses on symbolic reasoning, a task simple for humans but potentially difficult for language models. We demonstrate that chain-of-thought prompting not only allows language models to tackle symbolic reasoning problems that are tough with standard prompting but also enables length generalization to inference inputs that are longer than those seen in the few-shot examples.

\paragraph{Tasks.}
We consider the following two synthetic tasks.

\begin{itemize}[leftmargin=*,topsep=0pt]
    \itemsep0em \item \textbf{Last letter concatenation.} This task requires the model to join together the last letters of words in a name (e.g., \textit{``Amy Brown''} $\rightarrow$ \textit{``yn''}). 
    It is a more difficult variant of first letter concatenation, which language models can already handle without chain of thought.\footnote{We evaluated 10 common names using GPT-3 \texttt{davinci} and it answered all but one correctly.} We generate full names by randomly combining names from the top one-thousand first and last names based on name census data ({\small{\url{https://namecensus.com/}}}).
    \item \textbf{Coin flip.} This task requires the model to determine if a coin remains heads up after a series of people either flip or do not flip it (e.g., \textit{``A coin is heads up. Phoebe flips the coin. Osvaldo does not flip the coin. Is the coin still heads up?''} $\rightarrow$ \textit{``no''}).
\end{itemize}

Given that the creation of these symbolic reasoning tasks is clearly defined, for each task we examine both an \textit{in-domain} test set—where the examples have the same number of steps as the training/few-shot samples—and an \textit{out-of-domain} (OOD) test set, where the evaluation examples include more steps than those in the few-shot samples. In the last letter concatenation task, the model is only exposed to examples of names with two words, and is then tested on last letter concatenation involving names with 3 and 4 words.\footnote{For names longer than two words, we concatenate multiple first and last names together.} We apply a similar approach for the number of possible flips in the coin flip task. Our experimental setup utilizes the same methods and models as described in the previous two sections. We once more manually craft chains of thought for the few-shot exemplars for each task, as shown in \cref{fig:dataset-examples}.

\paragraph{Results.} 
The outcomes of these in-domain and OOD tests are presented in \cref{fig:symbolic-main} for \palm{}, with the results for \lamda{} reported in Appendix \cref{tab:all-lm-symbolic}. Using \palm{} 540B, chain-of-thought prompting achieves nearly perfect solve rates (note that standard prompting already solves the coin flip task with \palm{} 540, although not with \lamda{} 137B). It is important to highlight that these in-domain tests are considered ``toy tasks'' because the correct solution procedures are already supplied by the chains of thought in the few-shot examples; the model's role is simply to replicate the same steps using the new symbols in the test examples. Nevertheless, smaller models still fail— the capability to conduct abstract manipulations on unseen symbols across these three tasks only emerges at the scale of models with 100B parameters.

Regarding the OOD evaluations, standard prompting is ineffective for both tasks. However, with chain-of-thought prompting, language models exhibit upward scaling trends (although performance remains lower compared to the in-domain setting). Thus, chain-of-thought prompting enables length generalization beyond previously seen reasoning chains for sufficiently large language models.

\section{Discussion}\label{sec:discussion}

We have investigated chain-of-thought prompting as a straightforward approach to invoke multi-step reasoning capabilities in large language models. Initially, we observed that chain-of-thought prompting significantly enhances performance on arithmetic reasoning tasks, producing improvements far exceeding those from ablation studies and remaining consistent across various annotators, exemplars, and language models (\cref{sec:arithmetic-reasoning}). Subsequently, experiments involving commonsense reasoning highlighted the broadly applicable linguistic nature of chain-of-thought reasoning (\cref{sec:commonsense-reasoning}). Lastly, our findings show that for symbolic reasoning, chain-of-thought prompting aids out-of-distribution generalization to sequences of greater length (\cref{sec:symbolic-reasoning}). In all cases, chain-of-thought reasoning was triggered simply by prompting pretrained language models without any fine-tuning conducted during this study.

The rise of chain-of-thought reasoning correlated with model scale has been a consistent pattern \citep{wei2022emergent}. For many reasoning tasks where standard prompting yields flat scaling curves, chain-of-thought prompting produces sharply increasing scaling trajectories. This method seems to broaden the variety of tasks that large language models can effectively perform—in other words, our work emphasizes that standard prompting only establishes a lower bound on what large language models are capable of. This insight likely prompts more questions than it resolves—such as how much further reasoning ability may improve with additional scaling, and which other prompting strategies could further expand the scope of solvable tasks.

Regarding limitations, first, while chain-of-thought mimics human reasoning processes, this does not confirm that the neural network is genuinely “reasoning,” a question we leave open. Second, although the manual effort to augment exemplars with chains of thought is minor in few-shot settings, this annotation burden could be substantial for fine-tuning (though it might be mitigated through synthetic data generation or zero-shot generalization). \orange{Third, there is no assurance that the reasoning paths generated are correct, which means that both accurate and inaccurate answers can result; enhancing the factual reliability of language model outputs remains an important direction for future research \citep[][\textit{inter alia}]{rashkin2021measuring,ye2022unreliability,wiegreffe2021reframing}.} Finally, the fact that chain-of-thought reasoning only emerges at large model scales makes deployment costly in practical applications; further studies might investigate methods to elicit reasoning in smaller models.

\section{Related Work} This study draws inspiration from several research domains, which we elaborate on in a more comprehensive related work section (\cref{sec:extended-related-work}). Here, we focus on two key directions and their related literature that are particularly pertinent.

The first relevant avenue involves employing intermediate steps to address reasoning challenges. \cite{ling-etal-2017-program} were pioneers in utilizing natural language rationales to solve math word problems by breaking them down into a sequence of intermediate steps. Their approach stands in notable contrast to research that relies on formal languages for reasoning \citep{roy-2016-reasoning, chiang-chen-2019-semantically, amini-etal-2019-mathqa, chen2019neural}. Building on \citet{ling-etal-2017-program}, \citet{cobbe2021training} expanded the dataset size and finetuned a pretrained language model instead of training one from scratch. In the program synthesis field, \cite{nye2021show} utilize language models to predict the final outputs of Python programs by first sequentially predicting intermediate computational results, demonstrating that this stepwise prediction approach outperforms direct final output prediction.

Additionally, this paper is closely linked to the extensive recent research on prompting techniques. Since the rise of few-shot prompting popularized by \citet{brown2020language}, various strategies have enhanced models' prompting capabilities, including automated prompt learning \citep{lester-etal-2021-power} and providing task instructions to models \citep{wei2021finetuned,sanh2021multitask,ouyang2022training}. While these methods focus on improving or augmenting the input portion of prompts (such as prepending instructions), our approach diverges by augmenting the language model outputs with a chain of thought.

\section{Conclusions} We have investigated chain-of-thought prompting as a straightforward and widely applicable technique to boost reasoning abilities in language models. Our experiments spanning arithmetic, symbolic, and commonsense reasoning reveal that chain-of-thought reasoning emerges with model scale, enabling sufficiently large language models to tackle reasoning problems that otherwise display flat performance scaling. Expanding the variety of reasoning tasks language models can handle may encourage further exploration of language-driven reasoning methods.

\end{document}